\section{Conclusion and Future Work}

This work introduces \benchname, a unified benchmark for systematically evaluating memory-augmented robotic manipulation across four cognitive dimensions: temporal, spatial, object, and procedural memory. We further develop a family of vision-language-action (VLA) models and conducted controlled comparisons of symbolic, perceptual, and recurrent memory representations with multiple integration mechanisms. Results show that no single design consistently dominates: symbolic memory excels at counting and short-horizon reasoning, while perceptual memory is crucial for motion-centric and time-sensitive behaviors. 
% Overall, memory is best viewed as a task-dependent design choice rather than a one-size-fits-all solution.


Despite its breadth, \benchname\xspace focuses on tabletop manipulation with a fixed set of assets and evaluates a single pretrained backbone ($\pi_{0.5}$), leaving mobile manipulation and alternative architectures, such as memory-bank methods and additional VLA backbones, for future study. Our results further suggest that memory representations are complementary rather than exclusive, motivating unified frameworks that integrate multiple forms of memory. We hope \benchname\xspace serves as a foundation and lasting testbed for developing reliable, memory-augmented robotic generalist agents.